\section{Concrete Worked Examples}
\label{sec:appendix-worked-examples}
%% ====================================================================

This appendix presents fully-worked encodings of published MLIP studies
in the \MLIPsOntology{}. Each subsection takes one source paper and
extracts every fact that the ontology can express into Turtle listings,
ending with a discussion of what the source paper does \emph{not}
report---i.e., gaps that the ontology (or its aligned upper ontologies)
could record but that no paper is currently structured enough to
provide. These gaps are the practical motivation for the
\MLIPsOntology{}: making them visible is the first step towards filling
them in future publications.

\subsection{Kumar et al.\ (2025): MTP for TiCr$_2$-H Laves Phases}
\label{sec:appendix-example-kumar2025}

This subsection encodes the study by Kumar, K\"ormann, Grabowski, and
Ikeda~\cite{kumar2025ticr2h}. The authors train two Moment Tensor
Potentials (MTPs)---one per phase of the TiCr$_2$ Laves
compound---for hydrogen absorption over the concentration range
$0 < x \le 6$ in TiCr$_2$H$_x$. We encode the C15 (cubic) phase
artefacts; the C14 (hexagonal) artefacts follow the same pattern.

\paragraph{Material system.}
TiCr$_2$ is an intermetallic compound (composition AB$_2$,
$A=\mathrm{Ti}$, $B=\mathrm{Cr}$) and a Laves phase. The paper
considers two crystallographic forms; the example below uses the C15
cubic form. Crystal-structure attributes (space group, Pearson symbol,
Strukturbericht designation) are not first-class concepts of the
\MLIPsOntology{}; we record them as \texttt{materialClass} strings and
defer a structured \texttt{Phase} concept to future work
(\S\ref{sec:future-phase}).

\begin{lstlisting}[language=turtle,caption={Material system.}]
ex:mat-TiCr2 a mlips:MaterialSystem ;
    rdfs:label "TiCr2 (C15 Laves phase)" ;
    mlips:chemicalFormula "TiCr2" ;
    mlips:materialClass "Laves phase, C15 cubic (Fd-3m)" ;
    mlips:sameAsWikidata <http://www.wikidata.org/entity/Q15724720> .
\end{lstlisting}

\paragraph{DFT reference calculations.}
DFT calculations were performed in VASP using the projector
augmented-wave (PAW) method, the Perdew--Burke--Ernzerhof (PBE)
exchange--correlation functional, an energy cutoff of 400~eV, and a
$\Gamma$-centred $4 \times 4 \times 4$ k-point mesh with
Methfessel--Paxton smearing of width 0.1~eV. Calculations are
non-spin-polarised (the authors verified that magnetic moments do not
influence the predicted energies and forces). Each supercell is a
$2 \times 2 \times 2$ expansion of the primitive C15 cell, containing
48 metal atoms plus a variable number of hydrogen atoms.

\begin{lstlisting}[language=turtle,caption={DFT calculation and settings.}]
ex:dft-TiCr2H-c15 a mlips:DFTCalculation ;
    mlips:hasDFTSettings ex:dft-settings-TiCr2H-c15 .

ex:dft-settings-TiCr2H-c15 a mlips:DFTSettings ;
    mlips:usedDFTCode ex:lib-vasp ;
    mlips:xcFunctional "PBE" ;
    mlips:pseudopotentialType "PAW" ;
    mlips:energyCutoff 400.0 ;
    mlips:kPointMesh "4x4x4 Gamma-centred" .

ex:lib-vasp a mlips:Library ;
    rdfs:label "VASP" .
\end{lstlisting}

\paragraph{Training dataset.}
The C15 dataset reported in Step~1 of the active-learning pipeline
contains 1{,}019 atomic configurations (the C14 step-1 set has 3{,}766);
later steps (random hydrogen placement, ab~initio MD trajectories,
configurations sampled during basin-hopping Monte Carlo) further
extend the set. The dataset covers energies, forces, and stresses, and
the configurations span the chemical degree of freedom (variable $x$)
and the vibrational degree of freedom (MD snapshots). We record the
sampling strategy through the new \samplingStrategyR{} property
(\S\ref{sec:training-data-module}); concrete strategy individuals are
modelled as instances with \texttt{rdfs:label} until a controlled
vocabulary is developed in the journal extension.

\begin{lstlisting}[language=turtle,caption={Training dataset and sampling strategies.}]
ex:ds-TiCr2H-c15 a mlips:TrainingDataset ;
    rdfs:label "TiCr2-Hx training set, C15 phase, step 1" ;
    mlips:coversMaterial ex:mat-TiCr2 ;
    mlips:coversProperty mlips:Energy, mlips:Forces, mlips:Stresses ;
    mlips:datasetProvenance mlips:Published ;
    mlips:numConfigurations 1019 ;
    mlips:hasDFTCalculation ex:dft-TiCr2H-c15 ;
    mlips:samplingStrategy ex:chem-sampling, ex:vib-sampling,
                           ex:active-learning-sampling .

ex:chem-sampling a mlips:SamplingStrategy ;
    rdfs:label "Chemical sampling" ;
    rdfs:comment "Configurations across hydrogen concentrations 0 < x <= 6." .

ex:vib-sampling a mlips:SamplingStrategy ;
    rdfs:label "Vibrational sampling (DFT MD)" ;
    rdfs:comment "Snapshots from ab initio molecular dynamics at 500 K." .

ex:active-learning-sampling a mlips:SamplingStrategy ;
    rdfs:label "Active learning" ;
    rdfs:comment "Extrapolation-grade filtering with threshold gamma >~ 1." .
\end{lstlisting}

\paragraph{Method, hyperparameters, and implementation.}
The MTP method is fitted with the BFGS optimiser as implemented in the
MLIP-2 software. Of the four MTP hyperparameters declared by the
ontology, the paper reports the maximum complexity level $\levmax = 16$
explicitly; the cutoff radius $\Rcut$, minimum interatomic distance
$\Rmin$, and number of radial basis functions $\nRadialBasis$ are not
stated and are presumably defaults of MLIP-2. We record the explicit
setting and flag the missing values in
\S\ref{sec:appendix-example-kumar2025-missing}. The loss-function
weights ($1$, $0.01$, $0.001$ on energy, force, and stress times
volume) are method-level descriptors of the loss; we capture them in
the loss-function instance.

\begin{lstlisting}[language=turtle,caption={Method, training algorithm, and implementation.}]
ex:MTP a mlips:MLIPMethod ;
    rdfs:label "Moment Tensor Potential" ;
    mlips:hasFunctionalForm ex:mtp-functional-form ;
    mlips:hasLossFunction ex:mtp-loss-w1-001-0001 ;
    mlips:hasTrainingAlgorithm ex:bfgs ;
    mlips:hasHyperparameter ex:hp-rcut, ex:hp-rmin,
                            ex:hp-nq, ex:hp-levmax ;
    mlips:hasImplementation ex:mlip-package-v2 .

ex:mtp-functional-form a mlips:FunctionalForm ;
    rdfs:label "MTP moment-tensor polynomial" .

ex:mtp-loss-w1-001-0001 a mlips:LossFunction ;
    rdfs:label "Weighted MSE (energy 1, force 0.01, stress*V 0.001)" .

ex:bfgs a mls:Algorithm ;
    rdfs:label "BFGS" .

ex:mlip-package-v2 a mlips:Implementation ;
    mlips:implementedIn ex:lib-mlip ;
    mlips:version "2" .

ex:lib-mlip a mlips:Library ;
    rdfs:label "MLIP-2 package" ;
    schema:url <https://gitlab.com/ashapeev/mlip-2> .
\end{lstlisting}

\paragraph{MLIP run, hyperparameter settings, and trained model.}
The C15-MTP is the trained model produced by an MLIP run that applies
the MTP method on the C15 dataset. Only the explicit hyperparameter
setting (\texttt{lev\_max=16}) is recorded; the implicit settings are
discussed in \S\ref{sec:appendix-example-kumar2025-missing}.

\begin{lstlisting}[language=turtle,caption={MLIP run, hyperparameter settings, and trained model.}]
ex:run-mtp-TiCr2H-c15 a mlips:MLIPRun ;
    rdfs:label "MTP training run, C15 phase" ;
    mlips:appliesMethod ex:MTP ;
    mlips:runsOn ex:ds-TiCr2H-c15 ;
    mlips:hasHyperparameterSetting ex:setting-levmax-16 ;
    mlips:produces ex:model-mtp-TiCr2H-c15 .

ex:setting-levmax-16 a mlips:HyperparameterSetting ;
    mlips:forHyperparameter ex:hp-levmax ;
    mlips:settingValue "16" .

ex:model-mtp-TiCr2H-c15 a mlips:TrainedModel ;
    rdfs:label "C15-MTP for TiCr2-H" .
\end{lstlisting}

\paragraph{Benchmark study, result, and accuracy metrics.}
The paper reports the C15-MTP achieves a root-mean-square error of
$3.17$~meV/atom on energies (the C14-MTP reaches $2.81$~meV/atom).
Force errors are reported qualitatively as ``approximately within
3~meV/atom''; we record the per-property energy RMSE explicitly and
flag the missing force RMSE.

\begin{lstlisting}[language=turtle,caption={Benchmark study and result.}]
ex:study-kumar2025 a mlips:BenchmarkStudy ;
    rdfs:label "Kumar et al. (2025): MTP for TiCr2-H Laves phases" ;
    mlips:reportedIn ex:article-kumar2025 ;
    mlips:hasResult ex:result-mtp-ticr2h-c15 .

ex:article-kumar2025 a schema:ScholarlyArticle ;
    schema:sameAs <https://doi.org/10.1016/j.actamat.2025.121319> ;
    schema:datePublished "2025"^^xsd:gYear ;
    schema:name "Machine Learning Potentials for Hydrogen
                 Absorption in TiCr2 Laves Phases" .

ex:result-mtp-ticr2h-c15 a mlips:BenchmarkResult ;
    mlips:evaluatesModel ex:model-mtp-TiCr2H-c15 ;
    mlips:targetMaterial ex:mat-TiCr2 ;
    mlips:hasAccuracyMetric ex:metric-rmse-energy-c15 .

ex:metric-rmse-energy-c15 a mlips:AccuracyMetric ;
    mlips:metricType mlips:RMSE ;
    mlips:metricProperty mlips:EnergyProperty ;
    mlips:metricValue 3.17 ;
    mlips:hasUnit unit:MeV-PER-ATOM .
\end{lstlisting}

\subsubsection{What is missing from the paper}
\label{sec:appendix-example-kumar2025-missing}

The Kumar et al.\ paper is among the more methodologically explicit
MLIP papers, yet it leaves several pieces of metadata implicit or
absent. Each item below is expressible in the \MLIPsOntology{} (or in
an ontology we are aligned to) but is not provided by the source.

\paragraph{Hyperparameter values that defaulted silently.}
The MTP cutoff radius ($\Rcut$), minimum interatomic distance
($\Rmin$), and number of radial basis functions ($\nRadialBasis$) are
not stated. Reproducing the model requires either the MLIP-2 default
values at the time of training or the configuration files. The
ontology models these as \HyperparameterSettingC{} instances with no
recorded \texttt{settingValue}; consumers can detect this through a
SPARQL query and surface a warning.

\paragraph{Training-cost metadata.}
\trainingDurationR{}, \gpuHoursR{}, \trainingHardwareR{}, and
\peakMemoryR{} (data properties on \MLIPRunC{}) are all absent. CPU
hours, the number of DFT calls during active learning, and the cost
of the basin-hopping Monte Carlo loop would all be valuable for
comparing this work against alternative MLIP architectures (CQ9 in
the requirements).

\paragraph{Inference-cost metadata.}
\inferenceTimePerAtomR{}, \inferenceHardwareR{}, and \peakMemoryR{}
on the trained model are not reported. For a hydrogen-storage screening
application this matters: the practical value of the MTP comes from
its inference cost relative to DFT, which is alluded to but not
quantified.

\paragraph{Force-RMSE benchmark result.}
A second \BenchmarkResultC{} on the force property would parallel the
energy RMSE, but only the energy figure is given as a single number.
Our ontology trivially records this once it is available; the
underlying datum is missing from the paper.

\paragraph{Crystal-phase metadata as a first-class entity.}
The C15/C14 distinction is reduced to a string in
\texttt{materialClass}. Aligned ontologies (CMSO, EMMO) provide
crystal-structure, space-group, and Pearson-symbol concepts, but the
\MLIPsOntology{} does not currently expose these slots. Adding a
\dlConcept{Phase} class and aligning to CMSO's
\texttt{cmso:CrystalStructure} hierarchy is part of the planned
journal-extension work.
\label{sec:future-phase}

\paragraph{Step-by-step training provenance.}
The active-learning pipeline runs four distinct steps, each producing
an intermediate MTP and an extended training set. We collapse this
into a single \MLIPRunC{} on a single dataset; a faithful encoding
would chain four \MLIPRunC{} instances and four \TrainingDatasetC{}
versions linked through \texttt{prov:wasDerivedFrom}. The ontology
already supports this (datasets reuse PROV-O), but doing so was beyond
the scope of the original paper's reporting.

\paragraph{Reference-method and dataset deposit.}
Where the trained MTP and the training dataset live (DaRUS, Materials
Cloud, NOMAD, a personal repository) is not specified, so a
\datasetProvenanceR{} of \texttt{Published} is the only handle the
ontology has. Linking the dataset to a persistent identifier (DOI,
URN, or a Materials-Cloud entry IRI) is the kind of FAIR metadata
that the ontology is built to encode but that the paper does not yet
provide.

\paragraph{Note on non-DFT references.}
Kumar et al.\ used DFT throughout, so $\DFTCalculationC$ and
$\DFTSettingsC$ are the right slots. Had they instead used
CCSD(T) (or another wave-function method) for some configurations,
the encoding would replace $\DFTCalculationC$ with
$\WaveFunctionCalculationC$ and carry $\WaveFunctionSettingsC$ with
$\wfMethodR$~\texttt{= "CCSD(T)"}, $\basisSetR$~\texttt{= "cc-pVTZ"},
and $\frozenCoreR$~\texttt{= true}. The dataset-to-calculation link
$\hasReferenceCalculationR$ remains the same; the change is local to
the calculation subtype.

%% ====================================================================
\subsection{Extraction Protocol}
\label{sec:appendix-example-protocol}

The Kumar~et~al.\ subsection above is the reference for what an encoded
MLIP paper should look like. To replicate that quality across the
remaining subsections at scale---and to give a future agentic-AI tool
something concrete to follow---we fix a 12-question extraction
protocol. The protocol has two artefacts per paper: a canonical
Turtle file carrying the encoded data, and a subsection in this
appendix carrying the prose. The two are kept in sync by an automated
\emph{round-trip check} described at the end of this subsection.

\paragraph{Overview of the flow.}
For each paper we ask the same 12 questions in fixed order. Each
question (a) produces a Turtle fragment that is appended to the
paper's canonical file
\texttt{artifacts/kg/papers/\textit{paper-id}.ttl}, and (b) a short
commentary paragraph in the paper's subsection. Each of questions
Q1--Q11 has an associated CONSTRUCT query that, when run against the
paper's canonical file, regenerates the Turtle listing shown under
that question's commentary. The union of the 11 CONSTRUCT outputs is
required (modulo blank-node renaming and prefix declarations) to
reproduce the canonical file---this is the round-trip check. Q12 is
prose-only (gaps are not always reducible to triples) and is excluded
from the round-trip.

\paragraph{The 12 questions.}
Grouped into five phases:

\smallskip
\noindent\textbf{Phase A.\ Identification.}
\begin{itemize}
\item[Q1.] What are the bibliographic details of the source paper?
  Title, authors, year, venue, DOI.
  \\ \emph{Targets:} $\dlConcept{schema:ScholarlyArticle}$,
  $\BenchmarkStudyC$, $\reportedInR$.
\end{itemize}

\smallskip
\noindent\textbf{Phase B.\ Training data.}
\begin{itemize}
\item[Q2.] What material system(s) are studied? Chemical formula,
  material class (alloy, intermetallic, Laves phase, molecular,
  oxide), structural specifics (phase, space group, defect type),
  Wikidata link if available.
  \\ \emph{Targets:} $\MaterialSystemC$ and its data properties.
\item[Q3.] What method produced the reference data? DFT,
  wave-function (CCSD(T), MP2, CASPT2, \dots), AIMD, or experimental.
  \\ \emph{Targets:} $\ReferenceCalculationC$ and its concrete
  subtype.
\item[Q4.] What reference-method settings are reported? For DFT:
  $\xcFunctionalR$, $\kPointMeshR$, $\energyCutoffR$,
  $\pseudopotentialTypeR$. For wave-function: $\wfMethodR$,
  $\basisSetR$, $\frozenCoreR$. For both: software code via
  $\dlRole{usedReferenceCode}$ (or its DFT-specific sub-property).
  \\ \emph{Targets:} $\ReferenceSettingsC$ subtypes and $\LibraryC$.
\item[Q5.] How is the training dataset characterised? Number of
  configurations, covered physical properties (energies, forces,
  stresses, virials), provenance class, link to the reference
  calculation.
  \\ \emph{Targets:} $\TrainingDatasetC$, $\DatasetProvenanceC$,
  $\CoveredPropertyC$, $\hasReferenceCalculationR$.
\item[Q6.] What sampling strategies were used? Chemical, vibrational
  or MD-driven, active-learning, perturbation-based, exhaustive
  enumeration, or combinations.
  \\ \emph{Targets:} $\SamplingStrategyC$ instances and
  $\samplingStrategyR$.
\end{itemize}

\smallskip
\noindent\textbf{Phase C.\ Method, run, and model.}
\begin{itemize}
\item[Q7.] What MLIP method is used, and what are its components?
  Functional form, loss function, training algorithm (the ML-Schema
  alignment point), software implementation and version.
  \\ \emph{Targets:} $\dlConcept{MLIPMethod}$,
  $\dlConcept{FunctionalForm}$, $\dlConcept{LossFunction}$,
  $\hasTrainingAlgorithmR$, $\ImplementationC$, $\LibraryC$.
\item[Q8.] What hyperparameter settings does the paper explicitly
  report? Cutoff radius, MTP level, number of layers, learning rate,
  etc. Implicit defaults are recorded as
  \emph{not reported} but referenced in Q12.
  \\ \emph{Targets:} $\HyperparameterSettingC$ and
  $\forHyperparameterR$.
\item[Q9.] What is the structure of the run? Single training run vs.\
  multi-step active-learning pipeline; one trained model vs.\ a per-
  phase or per-fold ensemble.
  \\ \emph{Targets:} $\dlConcept{MLIPRun}$, $\TrainedModelC$,
  $\appliesMethodR$, $\runsOnR$, $\producesR$,
  $\hasHyperparameterSettingR$.
\end{itemize}

\smallskip
\noindent\textbf{Phase D.\ Evaluation.}
\begin{itemize}
\item[Q10.] What benchmark results does the paper report? Each
  accuracy metric is recorded with its type (RMSE, MAE, \dots), the
  property it measures (energy, forces, stresses, downstream
  properties), and its numerical value with units.
  \\ \emph{Targets:} $\BenchmarkResultC$, $\AccuracyMetricC$,
  $\MetricTypeC$, $\MetricPropertyC$, $\evaluatesModelR$,
  $\targetMaterialR$.
\item[Q11.] What computational-resource metadata is reported?
  Training duration, GPU hours, training hardware, peak memory;
  inference time per atom, inference hardware.
  \\ \emph{Targets:} data properties on $\dlConcept{MLIPRun}$ and
  $\TrainedModelC$.
\end{itemize}

\smallskip
\noindent\textbf{Phase E.\ Gaps.}
\begin{itemize}
\item[Q12.] What concepts that the ontology can express does the
  paper \emph{not} report? Free-form, prose-only. This question
  produces no Turtle and is excluded from the round-trip check.
\end{itemize}

\paragraph{Per-paper subsection template.}
A protocol-derived subsection contains 12 paragraphs corresponding
to Q1--Q12 (in order). Each paragraph for Q1--Q11 has the structure:
(i) a 2--4 sentence prose answer derived from the source paper;
(ii) a Turtle listing produced by the question's CONSTRUCT query;
and (iii) where the paper does not report a piece of data the
ontology can model, an explicit \emph{not reported} marker so that
Q12 has a complete inventory to summarise. Q12 is prose-only.

\paragraph{Files and naming.}
Per-paper data and queries live under \texttt{artifacts/kg/}:

\begin{itemize}
\item \texttt{papers/\textit{paper-id}.ttl}---canonical Turtle for
  one paper. Identifiers use a \texttt{firstauthor-year} convention
  (\texttt{kumar2025.ttl}, \texttt{qi2023.ttl}, \dots).
\item \texttt{queries/q01-bibliographic.rq} \dots
  \texttt{queries/q11-resources.rq}---one CONSTRUCT query per
  question Q1--Q11. The queries are paper-agnostic; they are
  parameterised only by the prefix declarations of the paper they
  are run against.
\item \texttt{check-roundtrip.sh \textit{paper-id}}---runs the 11
  CONSTRUCT queries against
  \texttt{papers/\textit{paper-id}.ttl}, takes their union, and
  diff-canonicalises against the canonical file. Exits non-zero on
  drift.
\item \texttt{build-listings.sh \textit{paper-id}}---runs the 11
  queries and emits Turtle listings ready for inclusion in the
  paper's subsection.
\end{itemize}

\paragraph{Round-trip check.}
At the end of every paper's encoding pass, we run
\texttt{check-roundtrip.sh}. The check parses the 11 query outputs
into N-triples (via \texttt{rapper -o ntriples}), sorts and dedups,
and diffs against the same canonicalisation of the paper's
\texttt{.ttl} file. Two outcomes are possible:

\begin{itemize}
\item \emph{Pass.} The 11 questions partition the paper's data; the
  subsection is a faithful presentation of the canonical file.
\item \emph{Fail.} A triple is in the canonical file but not in any
  CONSTRUCT output (encoded but never asked about), or vice versa.
  The protocol is then revised: either a question's scope is
  widened, or a new question is introduced and back-applied to all
  prior papers.
\end{itemize}

\paragraph{Limits and assumptions.}
Two assumptions worth flagging. First, the protocol assumes one
paper $\equiv$ one MLIP method $\times$ one material system. Papers
training multiple potentials or covering multiple systems
(e.g., Kumar et al.~2025 has separate C15 and C14 potentials)
instantiate the protocol per phase or per system, with the
appropriate identifiers (\texttt{kumar2025-c15},
\texttt{kumar2025-c14}). Second, papers occasionally report
ontology-relevant data the protocol does not capture (uncertainty
quantification, transferability claims, foundation-model fine-tuning
tags). When such data appears in three or more papers we treat that
as a signal to revise the protocol.

%% ====================================================================
